\section{Type does not imply occupation}

Let $V\simeq\mathbb C^5$ be the Higgs carrier and
$\Delta\in\operatorname{Sym}^2V$ the scalar-15 coordinate. After the body and
an occupied Higgs coordinate $H$ are frozen, the first local analytic action
has the form
\begin{equation}
\begin{split}
\Gamma_{15}[\Delta;H]={}&m_{\rm eff}^2(H)\lVert\Delta\rVert^2
+\lambda_1\lVert\Delta\rVert^4
+\lambda_2\operatorname{Tr}(\Delta^\dagger\Delta\Delta^\dagger\Delta)\\
&-2\operatorname{Re}\langle J_{15}(H,\psi),\Delta\rangle+\cdots,
\end{split}
\label{eq:localaction}
\end{equation}
where $m_{\rm eff}^2=M_\Delta^2+\kappa_{H\Delta}H^\dagger H$. An unbroken
$SU(5)$-covariant source obeys
$\operatorname{Hom}_{SU(5)}(15_H,\mathbf1)=0$, so it cannot provide a
constant nonsinglet load. A portal changes the Hessian but its first variation
vanishes at $\Delta=0$. If $m_{\rm eff}^2>0$, the quartic form is nonnegative
and no dual source or zero-excluding topology is present, the zero section is
the unique local minimum. Hence
\begin{equation}
\boxed{B_\tau^{\rm phys}+s_U+\text{current P3 law}
\not\Longrightarrow\Delta_*\ne0.}
\label{eq:occupationnogo}
\end{equation}

\begin{proposition}[Induced scalar-15 occupation]
If $H\ne0$, $K_{15}>0$ and the source owns a nonzero covariant incidence
$\mu_\Delta\langle(H\odot H)^\flat,\Delta\rangle+\mathrm{h.c.}$, then the
unique small-source branch satisfies
\begin{equation}
\boxed{\Delta_*=K_{15}^{-1}\bar\mu_\Delta(H\odot H)^\flat
+O(\lVert\mu_\Delta H\odot H\rVert^3).}
\label{eq:inducedoccupation}
\end{equation}
\end{proposition}
On the electroweak triplet component this gives
$v_\Delta\simeq\mu_\Delta v_H^2/M_{\Delta,\rm eff}^2$ and
$m_\nu\simeq Y_\Delta v_\Delta$. This proves neither $\mu_\Delta\ne0$ nor the
Yukawa map $Y_\Delta$.

\section{Typed seed edge and positive base realization}

The seed does not contain two Higgs particles. It contains an unordered
arity-two typed incidence
\begin{equation}
e_{HH\Delta}=(V_H,V_H;S_{15})/S_2,
\qquad H\mapsto H\odot H.
\label{eq:typededge}
\end{equation}
Let the enriched base realize it by an equivariant map
$C_{15}:\operatorname{Sym}^2V\to15_H$. Since both sides are the same
irreducible representation, Schur's lemma gives
$C_{15}=\mu_\Delta I_{15}$. For $K_{15}>0$, define
\begin{equation}
R_{15}=K_{15}^{1/2}\Delta-K_{15}^{-1/2}C_{15}(H\odot H).
\label{eq:matchingrow}
\end{equation}
Its once-counted positive action is
\begin{equation}
\begin{split}
\frac12\lVert R_{15}\rVert^2={}&
\frac12\langle\Delta,K_{15}\Delta\rangle
-\operatorname{Re}\langle C_{15}(H\odot H),\Delta\rangle\\
&+\frac12\langle C_{15}(H\odot H),
K_{15}^{-1}C_{15}(H\odot H)\rangle.
\end{split}
\label{eq:matchingaction}
\end{equation}
The middle term is the cubic incidence and the final term is its required
source-neutral completion, not a separately fitted Higgs interaction.
Variation gives $K_{15}\Delta_*=C_{15}(H\odot H)$.

Inside the frozen class requiring a renormalizable nonsinglet Majorana
mediator, one positive Higgs-dependent bosonic repair mode and one-carrier
source irredundancy, $15_H$ and Eq.~\eqref{eq:typededge} are the minimal common
carrier. Type-I and type-III fermionic mediators fail the bosonic one-carrier
requirement but remain alternatives outside this class. Neutrino mass alone
therefore does not select the branch.

\section{The cubic edge is not the complete RG packet}

The matching action contains stiffness, cubic incidence and a neutral $H^4$
completion, but no $H^2\Delta^2$ portal and no breaking descriptor that splits
the unified components. The weakest field-dependent stiffness is
\begin{equation}
K_{15}(H,\Sigma)=K_0+\kappa_{H\Delta}(H^\dagger H)I_{15}
+a_\Sigma\rho_{15}(\Sigma)+\cdots.
\end{equation}
Portal and breaking coefficients can vary while the cubic edge remains fixed;
this is an exact source-rank obstruction.

Assume local $SU(5)$ covariance, power-counting renormalizability, radiative
closure and no additional forbidding symmetry. A class-minimal support is
\begin{equation}
\begin{aligned}
I_{p1}&=(H^\dagger H)\operatorname{Tr}(\Delta\Delta^\dagger),&
I_{p2}&=H^\dagger\Delta\Delta^\dagger H,\\
I_b&=\langle\Delta,\rho_{15}(\Sigma)\Delta\rangle,&
I_{41}&=[\operatorname{Tr}(\Delta\Delta^\dagger)]^2,\\
&&I_{42}&=\operatorname{Tr}[(\Delta\Delta^\dagger)^2].
\end{aligned}
\label{eq:invariantsupport}
\end{equation}
This fixes operator directions, not the boundary vector
$c_I=(\kappa_1,\kappa_2,a_\Sigma,\lambda_{41},\lambda_{42})$. A common unit
cannot remove the remaining projective ratios. The finite selector
\begin{equation}
\mathcal L_I(c)=\tfrac12\langle c,G_Ic\rangle-\langle j_I,c\rangle,
\qquad c_I^*=G_I^{-1}j_I,
\end{equation}
is complete only after $G_I$ and $j_I$ are source-derived.

\section{Spectral-tensor boundary ray}

On $\mathcal H_{15}=\mathbb C\oplus V\oplus\bar V$, define
\begin{equation}
\mathcal D_{15}(H,\Delta,\Sigma)=
\begin{pmatrix}
0&H^\dagger&H^T\\ H&\Sigma&\Delta\\
H^*&\Delta^\dagger&\Sigma^T
\end{pmatrix}.
\label{eq:spectraloperator}
\end{equation}
It transforms by unitary conjugation under $SU(5)$. Its cubic moment contains
the breaking and charged incidences in one word,
\begin{equation}
\operatorname{Tr}\mathcal D_{15}^3
=6I_{H\Sigma}+6I_{HH\Delta}+3I_b+2I_{\Sigma3}.
\end{equation}
The canonical quartic completion
$\mathcal Q_4=\operatorname{Tr}\mathcal D_{15}^4
+(\operatorname{Tr}\mathcal D_{15}^2)^2$ gives at $\Sigma=0$
\begin{equation}
\boxed{\mathcal Q_4=24I_{H4}+16I_{p1}+8I_{p2}+4I_{41}+2I_{42}.}
\label{eq:spectralray}
\end{equation}
All directions in Eq.~\eqref{eq:invariantsupport} are present, and the portal
and scalar-quartic rays are each $2{:}1$. Equal weighting of single-trace and
tensor-square terms remains a completion premise, not a narrow-parent theorem.

\section{Common scale and metrological boundary}

The inherited common-length completion gives
\begin{equation}
\ell_H=\sqrt{48\pi}\,\ell_P,
\qquad
\boxed{M_*=\frac{\sqrt2\hbar c}{\ell_H}
=\frac{E_P}{\sqrt{24\pi}}
=\frac{\overline M_Pc^2}{\sqrt3}.}
\label{eq:commonscale}
\end{equation}
Using 2022 CODATA values \cite{CODATA2022} yields
$M_*=1.406035(16)\times10^{18}\,\mathrm{GeV}$. The Tau-specific statement is
the dimensionless ratio and absence of an independent scalar-15 length. The SI
number is a one-anchor calibration, not a prediction of $G$, $\hbar$ or $c$.

\section{Exact functional flow and its evaluation boundary}

Let $P$ be the positive background-covariant Laplace-type operator including
declared gauge-fixing and ghost blocks. The smooth spectral regulator
$R_k(P)=P/[\exp(P/k^2)-1]$ enters the exact effective-average-action flow
\cite{Wetterich1993}
\begin{equation}
\boxed{\partial_t\Gamma_k=\frac12\operatorname{STr}
[(\Gamma_k^{(2)}+R_k)^{-1}\partial_tR_k],\qquad t=\log k.}
\label{eq:wetterich}
\end{equation}
Gauge sectors must also satisfy the regulator-modified Slavnov--Taylor/BV
identity, whose insertion vanishes as $k\to0$. For a radiatively and
BRST-complete independent basis $\{\mathcal O_I\}$,
\begin{equation}
\boxed{\beta^I=(G^{-1})^{IJ}
\langle\mathcal O_J,\partial_t\Gamma_k\rangle,
\quad G_{IJ}=\langle\mathcal O_I,\mathcal O_J\rangle.}
\label{eq:betaprojection}
\end{equation}
These equations define an exact functional evolution and coordinate
projection after field content, scheme and matching are declared. They are
standard QFT machinery, not a Tau derivation of the Wetterich equation, and do
not evaluate the full component beta polynomials. Residual regulator
dependence in a truncation is an uncertainty diagnostic.

For one active complex $15_H$, representation theory fixes
$\Delta b^{15_H}=(7/6,7/6,7/6)$. This changes common running but not relative
one-loop gauge slopes. A unique trajectory still requires portals, scalar
quartics, Yukawa/cubic data, component thresholds and breaking-scale matching.
Equal-upstream choices of undetermined data can produce different
stationarity verdicts, so a chosen trajectory is not a Tau prediction.
